\documentclass[10pt]{article}
\usepackage[T1]{fontenc}
\usepackage[utf8]{inputenc}
\usepackage[french]{babel}
\usepackage{amsmath,amssymb,mathrsfs,mathpazo}
\usepackage{geometry}
\geometry{a4paper,margin=2cm}
\usepackage{graphicx,array,tabularx,tabularray,colortbl,multirow,multicol,booktabs,enumitem,pifont,ulem,fancybox,diagbox,cancel}
\usepackage{pstricks,pstricks-add,pst-plot,pst-tree,pst-node,pst-text,pst-func,pst-math,pst-bspline,pst-3dplot,pst-eucl,pst-solides3d}
\let\pscalculate\undefined
\usepackage{pst-bezier}
\usepackage{tikz,pgfplots,tkz-tab}
\usetikzlibrary{arrows,arrows.meta,patterns,calc,positioning,trees,decorations.markings,shapes.geometric,shapes.misc,shapes.arrows,intersections,decorations.pathmorphing,decorations.pathreplacing,backgrounds,shadows,babel}
\pgfplotsset{compat=1.18}
\usepgfplotslibrary{fillbetween}
\usepackage[np]{numprint}
\usepackage{listings,makecell,pas-tableur,twoopt}
\let\barre\undefined
\usepackage[pstricks]{tablvar}
\definecolor{darkgreen}{RGB}{0,100,0}
\providecommand{\red}{\color{red}}
\providecommand{\blue}{\color{blue}}
\providecommand{\green}{\color{green}}
\providecommand{\black}{\color{black}}
\providecommand{\white}{\color{white}}
\providecommand{\gray}{\color{gray}}
\providecommand{\cyan}{\color{cyan}}
\providecommand{\magenta}{\color{magenta}}
\providecommand{\yellow}{\color{yellow}}
\providecommand{\violet}{\color{violet}}
\providecommand{\orange}{\color{orange}}
\providecommand{\euro}{\texteuro}
\providecommand{\EUR}{\texteuro}
\providecommand{\ieme}{\textsuperscript{e}}
\providecommand{\degres}{\ensuremath{^\circ}}
\providecommand{\degree}{\ensuremath{^\circ}}
\providecommand{\vv}[1]{\overrightarrow{#1}}
\providecommand{\Coord}[1]{#1}
\providecommand{\decosix}{}
\providecommand{\starredbullet}{\ensuremath{\bullet}}
\pagestyle{empty}
\setlength{\parindent}{0pt}
\usepackage{adjustbox}\DeclareUnicodeCharacter{2236}{:}\DeclareUnicodeCharacter{2190}{\ensuremath{\leftarrow}}
\providecommand{\R}{\mathbb{R}}
\providecommand{\N}{\mathbb{N}}
\providecommand{\Z}{\mathbb{Z}}
\providecommand{\Q}{\mathbb{Q}}
\providecommand{\D}{\mathbb{D}}
\providecommand{\C}{\mathbb{C}}
\providecommand{\e}{{\mathrm{e}}}
\providecommand{\vect}[1]{\overrightarrow{#1}}
\providecommand{\vectt}[1]{\overrightarrow{\mathrm{#1}}}
\providecommand{\vvt}[1]{\overrightarrow{\mathrm{#1}}}
\providecommand{\np}[2][]{{#2~#1}}
\providecommand{\vv}[1]{{\overrightarrow{#1}}}
\renewcommand{\d}{{\mathrm{d}}}
\providecommand{\degres}{{^\circ}}
\providecommand{\degre}{{^\circ}}
\providecommand{\blue}{\color{blue}}
\providecommand{\red}{\color{red}}
\providecommand{\enskip}{\quad}
\providecommand{\llbracket}{\lbrack\!\lbrack}
\providecommand{\rrbracket}{\rbrack\!\rbrack}
\providecommand{\euro}{{\text{€}}}
\providecommand{\no}{{\text{n°}}}
\providecommand{\interval}[2]{{[#1;#2]}}
\providecommand{\intervFO}[2]{{[#1;#2[}}
\providecommand{\dotfill}{\ldots}

\providecommand{\ds}{\displaystyle}
\providecommand{\textding}[1]{#1}
\providecommand{\Oij}{\(\mathrm{O};\vec{\imath},\vec{\jmath}\)}
\providecommand{\Oijk}{\(\mathrm{O};\vec{\imath},\vec{\jmath},\vec{k}\)}
\providecommand{\Ouv}{\(\mathrm{O};\vec{u},\vec{v}\)}
\providecommand{\intervFF}[2]{\ensuremath{[#1;#2]}}
\renewcommand{\euro}{\ifmmode\text{\texteuro}\else\texteuro\fi}
\DeclareUnicodeCharacter{00B0}{\ifmmode^{\circ}\else\textdegree\fi}
\DeclareUnicodeCharacter{20AC}{\ifmmode\text{\texteuro}\else\texteuro\fi}
% Énoncés : Éducation nationale. Transcription et corrigés : APMEP, Denis Vergès et contributeurs. Découpage et classement : CoopMaths, Sébastien Lozano, Éric Elter et contributeurs. Adaptation de la présentation : Maths Explorer.
% CC BY-SA : https://coopmaths.fr/www/about/
% https://coopmaths.fr/alea/static/e3c/2020/tex/e3c_2020_00_sujet56_1.tex
\begin{document}
\noindent\textit{APMEP / CoopMaths. CC BY-SA. Présentation adaptée par Maths Explorer.}\par\medskip
\medskip

\emph{Ce QCM comprend $5$ questions indépendantes.}

\emph{Pour chacune d'elles, une seule des réponses proposées est exacte.}

\emph{Indiquer pour chaque question sur la copie la lettre correspondant à la réponse
choisie. Aucune justification n'est demandée.}

\emph{Chaque réponse correcte rapporte 1 point. Une réponse incorrecte ou une absence de réponse n'apporte ni ne retire de point.}

\medskip

\textbf{QUESTION 1}

\medskip

\begin{tabularx}{\linewidth}{*{4}{X}}
\textbf{a.~~} Si le discriminant
d'un polynôme
du second degré
est strictement
positif, alors ce
polynôme admet
2 racines
positives.&
\textbf{b.~~} Si le discriminant
d'un polynôme
du second degré
est strictement
négatif, alors ce
polynôme admet
2 racines
négatives.&
\textbf{c.~~} Si un polynôme
du second degré
est toujours
strictement
positif, alors ce
polynôme
n'admet pas de
racine.&
\textbf{d.~~}Si le discriminant
d'un polynôme
du second degré
est nul, alors ce
polynôme admet
le nombre 0 pour
racine.
\end{tabularx}

\medskip

\textbf{QUESTION 2}

\medskip

\begin{tabularx}{\linewidth}{*{4}{X}}
\begin{minipage}[t]{3.5cm}\textbf{a.~~}
L'équation

 $\cos x =-\frac{1}{2}$
admet 2 solutions dans l'intervalle
 $\left]-\frac{\pi}{2};\frac{\pi}{2}\right]$.
\end{minipage}&
\begin{minipage}[t]{3.5cm}\textbf{b.~~}
 L'équation
 
$\cos x= -\frac{1}{2}$
admet 1
solution dans
l'intervalle $[0 ; \pi[$.
\end{minipage}&
 \begin{minipage}[t]{3.5cm}\textbf{c.~~}
L'équation

$\sin x = -\frac{1}{2}$ admet 1
solution dans
l'intervalle \mbox{$[0 ; \pi[$}.
\end{minipage}&
 \begin{minipage}[t]{3.5cm}\textbf{d.~~}
\L'équation

$\sin x= -\frac{1}{2}$ admet 2
solutions dans
l'intervalle $\left]-\frac{\pi}{2} ;\frac{\pi}{2}\right]$.
\end{minipage}
\end{tabularx}

\medskip

\textbf{QUESTION 3}

\medskip

La courbe représentative d'une fonction $f$, définie et dérivable sur l'ensemble
des nombres réels, est donnée ci-dessous avec ses tangentes, aux points A
et B d'abscisses respectives 2 et 4. On note $f'$ la fonction dérivée de $f$.

\begin{center}
\psset{unit=0.5cm,labelFontSize=\scriptstyle,labelsep=0.1pt,showorigin=false}
\begin{pspicture}(-3.5,-5.8)(9,7)
\multido{\n=-3+1}{13}{\psline[linewidth=0.75pt,linecolor=lightgray](\n,-4.5)(\n,6.2)}
\multido{\n=-4+1}{11}{\psline[linewidth=0.75pt,linecolor=lightgray](-3.5,\n)(9.2,\n)}
\psaxes[linewidth=0.95pt,]{->}(0,0)(-3.5,-4.5)(9.5,6.4)
\psdots[dotstyle=+,dotscale =1.4,dotangle=45](2,-2) \psdots[dotstyle=+,dotscale =1.4,dotangle=45](4,-1)
\uput[dl](3.8,-1){B} \uput[dr](2,-2.1){A} 
\psplot[linewidth=1pt,linecolor=blue,plotpoints=5000]{-2.7}{5.7}{ x x x 0.125 neg mul 0.75 add mul 0.5 sub mul 3 sub }
\psplot[linewidth=1pt,linecolor=orange,plotpoints=3000]{-1}{9}{x  4 sub}
\psplot[linewidth=1pt,linecolor=red,plotpoints=3000]{-3}{9}{x 0.5 neg mul 1  add}
\end{pspicture}
\end{center}

\begin{tabularx}{\linewidth}{*{4}{X}}
\textbf{a.~~} $f(0) = 1$ &\textbf{b.~~} $f'(2) = 1$&\textbf{c.~~}$f'(2) = -2 $& \textbf{d.~~} $f'(4) = 0,5$.
\end{tabularx}
\medskip

\textbf{QUESTION 4}

\medskip

On considère la fonction $g$ définie sur l'ensemble des nombres réels $\R$ par :
$g(x)= x^3 - \np{0,0012}x + 1$.

\medskip

\begin{tabularx}{\linewidth}{*{4}{X}}
\textbf{a.~~} $g$ est strictement
croissante sur~$\R$.&
\textbf{b.~~} $g$ est croissante
sur~$R$.&
\textbf{c.~~} $g$ est constante
sur l'intervalle
\mbox{$[- 0,02~;~0,02]$}.&
\textbf{d.~~} $g$ est
décroissante sur
l'intervalle
\mbox{$[- 0,02~;~0,02]$}.
\end{tabularx}

\medskip

\textbf{QUESTION 5}

\medskip
\begin{tabularx}{\linewidth}{*{4}{X}}

\textbf{a.~~} L'équation
\mbox{$(\e^x)^2 = 1$} 
admet deux solutions
dans $\R$.

&\textbf{b.~~} L'ensemble de
définition de la
fonction
exponentielle est
\mbox{$]0~;~+ \infty[$}.

&\textbf{c.~~} La fonction
dérivée de la
fonction
$x\mapsto \e^{-x}$ est la
fonction $x \mapsto \e^{-x}$.

& \textbf{d.~~} L'ensemble de
définition de la
fonction
exponentielle est
$\R$.
\end{tabularx}

\vspace{0,5cm}

\end{document}
